% =====================================================================
%  Architekturübersicht BalanceOS: vier Schichten (Python-Pakete) mit
%  Abhängigkeitsrichtung und Artefakten (SQLite-Datei, Exportdateien).
%  Abgelesen aus app/src/balanceos (Code-Freeze Phase 2).
%  Export: docs/diagramme/export/architektur_uebersicht.pdf
% =====================================================================
\documentclass[tikz,border=4pt]{standalone}
\usepackage{iu-diagramm}

\tikzset{
  schicht/.style={draw, fill=iuFlaeche, rounded corners=3pt, align=left, inner sep=5pt,
                  text width=#1, anchor=north west},
  schicht/.default=14.6cm,
  artefakt/.style={draw, fill=white, align=center, inner sep=4pt, text width=3.4cm,
                   font=\fontsize{8pt}{9.5pt}\selectfont,
                   path picture={\draw[iuLinie] ([xshift=-2.2mm]path picture bounding box.north east)
                     -- ([yshift=-2.2mm]path picture bounding box.north east);}},
  abh/.style={dashed, -{Straight Barb[length=2.2mm]}},
  klein/.style={font=\fontsize{8pt}{9.5pt}\selectfont},
}
\newcommand{\kopf}{\fontsize{9.5pt}{11pt}\selectfont\bfseries}
\newcommand{\klein}{\fontsize{8pt}{9.5pt}\selectfont}

\begin{document}
\begin{tikzpicture}

\node[schicht] (pr) at (0,0) {%
  {\kopf «layer» Präsentation \hfill\normalfont\ttfamily praesentation/}\\[2pt]
  \klein Hauptfenster mit Seitenleiste, sieben Seiten (Dashboard, Woche, Aktivitäten, Statistik,
  Reflexion, Kompass, Einstellungen), Dialoge, eigene Widgets (Wochenkalender, Soll-Ist-Balken,
  Liniendiagramm), \texttt{Zustand} mit den Qt-Signalen \texttt{daten\_geaendert} und
  \texttt{woche\_geaendert}\\[2pt]
  \klein\itshape nutzt PySide6 (Qt Widgets)};

\node[schicht] (an) at (0,-2.75) {%
  {\kopf «layer» Anwendung \hfill\normalfont\ttfamily anwendung/}\\[2pt]
  \klein \texttt{AppKontext} (Verdrahtung), Services: \texttt{StammdatenService},
  \texttt{AktivitaetService}, \texttt{AuswertungService}, \texttt{ReflexionService},
  \texttt{ExportService}, \texttt{DatenService}; Demodaten; Fehlerklassen
  \texttt{ValidierungsFehler}, \texttt{RegelVerletzung}};

\node[schicht=6.3cm] (do) at (0,-5.15) {%
  {\kopf «layer» Domäne \hfill\normalfont\ttfamily domaene/}\\[2pt]
  \klein Datenklassen und Aufzählungen (\texttt{modell}, \texttt{typen}), \texttt{validierung}
  (GR01--GR30), \texttt{bewertung} (Score), \texttt{regeln} (R01--R10), \texttt{serien},
  \texttt{kalender} (ISO-Woche), \texttt{who5}\\[2pt]
  \klein\itshape nur Python-Standardbibliothek, ohne Qt und \mbox{SQLite}};

\node[schicht=6.3cm] (pe) at (8.3,-5.15) {%
  {\kopf «layer» Persistenz \hfill\normalfont\ttfamily persistenz/}\\[2pt]
  \klein \texttt{Datenbank} (Verbindung, Schema, Transaktionen), Repositories je Aggregat,
  Stammdaten (zehn Standardbereiche, Regeln, Kompass-Modelle)\\[2pt]
  \klein\itshape nutzt \texttt{sqlite3} (Standardbibliothek)};

\node[artefakt, anchor=north] (db) at ([yshift=-0.9cm]pe.south) {«file»\\\texttt{balanceos.db}\\im Benutzerprofil (\%APPDATA\%)};
\node[artefakt, anchor=west] (ex) at ([xshift=1.5cm]an.east) {«file»\\CSV- und JSON-Export\\(von der Nutzer:in gewählter Ordner)};

% Abhängigkeiten (gestrichelt, Pfeil zur genutzten Schicht)
\draw[abh] ([xshift=3cm]pr.south west) -- node[klein, right] {nutzt Services} ([xshift=3cm]an.north west);
\draw[abh] (pr.west) -- ++(-0.45,0) |- node[klein, pos=0.25, left, align=right] {Datenklassen,\\Feld-\\prüfung} (do.west);
\draw[abh] ([xshift=3cm]an.south west) -- node[klein, right] {Fachlogik} ([xshift=3cm]do.north west);
\draw[abh] ([xshift=3cm]pe.north west |- an.south) -- node[klein, right] {Repositories} ([xshift=3cm]pe.north west);
\draw[abh] (pe.west) -- node[klein, above] {Daten-} node[klein, below] {klassen} (do.east);
\draw[abh] (pe.south) -- node[klein, right] {SQL} (db.north);
\draw[abh] (an.east) -- node[klein, above] {schreibt} (ex.west);

\node[klein, anchor=north west, align=left, text width=14.6cm] at ([yshift=-0.35cm]do.south west |- db.south) {%
  Gestrichelte Pfeile: Abhängigkeit (Import) in Pfeilrichtung. Die Präsentation greift über die
  Anwendungsschicht auf die Daten zu und nutzt aus der Domäne nur Datenklassen und Prüffunktionen
  (lockere Schichtung). Einzige Ausnahme ist der Programmstart \texttt{praesentation/app.py}, der
  den Speicherort der Datenbank aus der Persistenz erfragt.};

\end{tikzpicture}
\end{document}
